\section{Original GSL/cGSL Results}\label{app:original_gsl}

This appendix presents the original Graph Structure Learning (GSL) and cyclic GSL (cGSL) results that motivated the multi-lag investigation. Method naming in this appendix follows the original submission; all tables in this appendix are \textbf{historical single-seed results} produced under the original study's protocol and pipeline. They are retained for transparency and because the GSL/cGSL backbone asymmetry they document remains scientifically informative; they are \emph{not} current canonical results and are not directly comparable in RMSE magnitude to the main-text tables, which use the revised training pipeline. The main-text re-establishment of this baseline under the revised protocol is given in Section~\ref{sec:gsl_baseline}.

\subsection{Method}

The original GSL approach applies DAGMA to contemporaneous traffic observations $X \in \mathbb{R}^{T \times N}$, where each row is a single-timestep snapshot of all $N$ sensor readings. This produces a single $N \times N$ weight matrix $W$. The GSL variant uses the binary thresholded DAG adjacency: $A_{GSL} = \mathbf{1}(|W| > \omega)$. The cGSL variant symmetrizes: $A_{cGSL} = A_{GSL} + A_{GSL}^T$.

Two protocol details, established during the revision by re-deriving the construction from the original data-loading code, complete the historical description. First, in the as-submitted pipeline DAGMA was not fitted once per horizon: for prediction horizon PH, the loader fitted $K = \mathrm{PH}$ separate DAGMA models on stride-$K$ offset-stratified subsamples of the training windows ($X = \mathrm{data}[i::K]$ for $i = 0, \ldots, K{-}1$) and used the \emph{union} of the $K$ resulting supports. The resulting as-used Los-loop graphs had 28/32/33/39 edges for PH=1--4 (the union of several DAGs need not itself be acyclic). Second, the library's absolute-magnitude threshold ($\omega = 0.3$) was applied inside each fit, and the historical union rule additionally retained positive coefficients only.

The revised protocol (Section~\ref{sec:gsl_baseline}) instead fits one DAGMA model per horizon on the offset-0 subsample, yielding a genuine single 28-edge DAG per horizon, and retains both coefficient signs after the library threshold. Its slice-0 fits reproduce the historical slice-0 graph supports exactly at PH=1, 2, and 4 (28/28 edges; 26/28 at PH=3), with shared-edge weights agreeing to within approximately 0.016; the revised per-horizon graphs are strict subsets of the historical unions. The historical tables below remain unchanged single-seed results of the original pipeline.

\subsection{GCN Results}

Table~\ref{tab:gcn_gsl_combined} shows the GCN results with physical graph, GSL, and cGSL.

\begin{table}[!t]
\centering
\caption{GCN performance with physical graph, GSL, and cGSL on both datasets (PH=1--4).}
\label{tab:gcn_gsl_combined}
\small
\begin{tabular}{cl|cc|cc}
\toprule
& & \multicolumn{2}{c}{SZ-taxi} & \multicolumn{2}{c}{Los-loop} \\
PH & Method & RMSE & MAE & RMSE & MAE \\
\midrule
& GCN & 5.958 & 4.409 & 7.724 & 5.555 \\
1 & GCN-GSL & 4.886 & 3.161 & 7.527 & 5.065 \\
& \textbf{GCN-cGSL} & \textbf{4.648} & \textbf{3.078} & \textbf{5.440} & \textbf{3.452} \\
\midrule
& GCN & 5.983 & 4.437 & 7.940 & 5.692 \\
2 & GCN-GSL & 4.904 & 3.173 & 7.867 & 5.268 \\
& \textbf{GCN-cGSL} & \textbf{4.672} & \textbf{3.051} & \textbf{5.806} & \textbf{3.613} \\
\midrule
& GCN & 5.991 & 4.438 & 8.102 & 5.782 \\
3 & GCN-GSL & 4.958 & 3.223 & 8.073 & 5.453 \\
& \textbf{GCN-cGSL} & \textbf{4.712} & \textbf{3.122} & \textbf{6.171} & \textbf{3.821} \\
\midrule
& GCN & 6.002 & 4.450 & 8.285 & 5.876 \\
4 & GCN-GSL & 4.933 & 3.199 & 9.067 & 6.085 \\
& \textbf{GCN-cGSL} & \textbf{4.726} & \textbf{3.107} & \textbf{6.745} & \textbf{4.097} \\
\bottomrule
\end{tabular}
\end{table}

GCN-cGSL (symmetrized cyclic graph) consistently outperforms both standard GCN and GCN-GSL, with mean RMSE improvements over the physical graph of 21.6\% on SZ-taxi and 24.7\% on Los-loop (multi-horizon means over PH=1--4, computed from Table~\ref{tab:gcn_gsl_combined}). This suggests that the purely spatial GCN benefits from bidirectional (cyclic) graph connections.

\subsection{TGCN Results}

Table~\ref{tab:tgcn_gsl_combined} shows the TGCN results.

\begin{table}[!t]
\centering
\caption{TGCN performance with physical graph, GSL, and cGSL on both datasets (PH=1--4).}
\label{tab:tgcn_gsl_combined}
\small
\begin{tabular}{cl|cc|cc}
\toprule
& & \multicolumn{2}{c}{SZ-taxi} & \multicolumn{2}{c}{Los-loop} \\
PH & Method & RMSE & MAE & RMSE & MAE \\
\midrule
& TGCN & 4.866 & 3.606 & 6.588 & 4.573 \\
1 & \textbf{TGCN-GSL} & \textbf{4.214} & \textbf{2.797} & \textbf{4.818} & \textbf{2.980} \\
& TGCN-cGSL & 4.821 & 3.537 & 6.550 & 4.553 \\
\midrule
& TGCN & 4.506 & 3.213 & 6.960 & 4.838 \\
2 & \textbf{TGCN-GSL} & \textbf{4.239} & \textbf{2.801} & \textbf{5.400} & \textbf{3.410} \\
& TGCN-cGSL & 4.534 & 3.276 & 6.915 & 4.830 \\
\midrule
& TGCN & 4.685 & 3.416 & 7.361 & 5.053 \\
3 & \textbf{TGCN-GSL} & \textbf{4.344} & \textbf{3.062} & \textbf{5.846} & \textbf{3.467} \\
& TGCN-cGSL & 4.630 & 3.334 & 7.331 & 5.054 \\
\midrule
& TGCN & 4.934 & 3.638 & 7.568 & 5.166 \\
4 & \textbf{TGCN-GSL} & \textbf{4.366} & \textbf{2.885} & \textbf{6.257} & \textbf{3.805} \\
& TGCN-cGSL & 4.774 & 3.487 & 7.539 & 5.172 \\
\bottomrule
\end{tabular}
\end{table}

TGCN-GSL (acyclic DAG graph) outperforms both standard TGCN and TGCN-cGSL, achieving a mean 21.8\% RMSE improvement over the physical graph on Los-loop (multi-horizon mean over PH=1--4, single seed, original pipeline; per-horizon maxima are higher, e.g., PH=1: 6.588 $\rightarrow$ 4.818, a 26.9\% reduction). The acyclic structure aligns well with T-GCN's temporal modeling through the GRU component. Under the revised pipeline this comparison is reproduced at five seeds with a 21.7\% mean improvement (25.5\% at PH=1; Section~\ref{sec:gsl_baseline}), which independently validates the magnitude of the original claim even though the absolute RMSE values of the two pipelines are not directly comparable.

\subsection{Key Observation}

The asymmetry between GCN (cyclic preferred) and T-GCN (acyclic preferred) motivated the investigation into what the learned graph actually represents. The multi-lag formulation (main text) resolves this by making the temporal structure explicit rather than relying on post-hoc interpretation of a single ambiguous graph. The reproduction analysis additionally showed that this original single-graph construction was contemporaneous (Section~\ref{sec:gsl_baseline}): the DAGs above encode present-observation functional dependencies, and the temporal interpretation is supplied by the explicitly lagged multi-lag construction of Section~\ref{sec:multilag}.
